% !PRINTED 4
\weberpage{4}

%% [running head, not body text -- commented out 2026-08-07] Page furniture (running
%% head, folio, signature) is excluded from this edition by standing policy.
%% \noindent{\small 4 \hfill Erster Abschnitt. \hfill \S.\,1.}

%<<CONT
and from this it follows, by the mode of inference of complete induction, that one always arrives at the same result when, in an arbitrary sequence of elements from $P$ in finite number, $a, b, c, d \ldots$, one first compounds two adjacent elements, then again two adjacent ones, and so on, until the whole sequence is reduced to a single element, which is denoted by $abcd \ldots$. Thus, for example:
\[
\begin{aligned}
abcd &= (ab)cd = [(ab)c]d = (ab)(cd)\\
&= a(bc)d = [a(bc)]d = a[(bc)d]\\
&= ab(cd) = (ab)(cd) = a[b(c.d)]^{1)}. % [sic: the source prints a dot between c and d]
\end{aligned}
\]

\noindent 3.\ \gesperrt{It is presupposed. that, if $ab = ab'$ or $ab = a'b$, then necessarily $b = b'$ in the first case, and $a = a'$ in the second.} % [sic: the source prints a full stop after vorausgesetzt]

If $P$ comprises a finite number of elements, the group is called a \gesperrt{finite one}, and the number of its elements its \gesperrt{degree}.

For finite groups there results from 1., 2., 3. the \gesperrt{consequence:}

\noindent 4.\ \gesperrt{If, of the three elements $a, b, c$ from $P$, any two are given, then the third can always be determined, and in only one way, so that}
\[
ab = c
\]
\gesperrt{holds.}

For if $a, b$ are the given elements, the assertion 4. coincides with 1. But if $a$ and $c$ are given, let the second component $b$ in the compositum $ab$ run through the whole system $P$, whose degree may be $= n$. Then by 1. and 3. one obtains in $ab$ nothing but distinct elements of $P$, and since their number is $= n$, all the elements of $P$, hence also $c$, must occur among them. One reasons in the same way when $b$ and $c$ are given, by letting $a$ run through the whole system $P$.

For infinite groups one can no longer reason in this way. \gesperrt{For infinite groups, therefore, the property 4. is further taken up into the definition as a requirement.}

\noindent\rule{3cm}{0.4pt}

\noindent{\footnotesize $^{1)}$ Cf.\ \gesperrt{Dirichlet-Dedekind}, Vorlesungen \"uber Zahlentheorie, \S.\,2.}
